\section{Planes in Three-Dimensional Space}

A line in $\mathbb{R}^3$ is determined by one point and one direction. A plane
requires one point and two independent directions. Equivalently, it may be
described by one point and a direction perpendicular to the whole plane.

\subsection{A Plane Through a Point with Two Direction Vectors}

Let
\[
P_0=(x_0,y_0,z_0)
\]
be a point and let
\[
\mathbf{u}=(u_1,u_2,u_3),
\qquad
\mathbf{v}=(v_1,v_2,v_3)
\]
be non-parallel vectors.

\begin{definitionbox}
The plane through $P_0$ generated by $\mathbf{u}$ and $\mathbf{v}$ is
\[
\Pi
=
\left\{
P_0+s\mathbf{u}+t\mathbf{v}:s,t\in\mathbb{R}
\right\}.
\]
In coordinates,
\[
\begin{cases}
x=x_0+s u_1+t v_1,\\
y=y_0+s u_2+t v_2,\\
z=z_0+s u_3+t v_3.
\end{cases}
\]
\end{definitionbox}

The parameters $s$ and $t$ allow independent motion in two directions. Because
$\mathbf{u}$ and $\mathbf{v}$ are not parallel, these motions fill a plane rather
than a single line.

\begin{center}
\begin{tikzpicture}[scale=0.9]
    \draw[axis] (-0.5,0) -- (5.4,0) node[right] {$x$};
    \draw[axis] (0,-0.4) -- (0,4.4) node[above] {$z$};
    \draw[axis] (0,0) -- (-2.0,-1.5) node[below left] {$y$};

    \filldraw[fill=geometricSubsetColor!12,draw=geometricSubsetColor,thick]
        (-1.2,0.8) -- (3.8,0.8) -- (5.0,3.1) -- (0.0,3.1) -- cycle;
    \coordinate (P) at (1.5,1.8);
    \fill[geometricSubsetColor] (P) circle (2.2pt) node[below left] {$P_0$};
    \draw[->,thick,geometricSubsetColor] (P) -- ++(2.0,0) node[right] {$\mathbf{u}$};
    \draw[->,thick,geometricSubsetColor] (P) -- ++(0.9,1.1) node[above] {$\mathbf{v}$};
    \node[subset label] at (4.3,2.8) {$\Pi$};
\end{tikzpicture}
\end{center}

\subsection{The Normal Vector and the Cartesian Equation}

A nonzero vector
\[
\mathbf{n}=(a,b,c)
\]
is called a normal vector to a plane when it is perpendicular to every direction
contained in that plane.

For a point $P=(x,y,z)$ on the plane, the displacement from $P_0$ to $P$ is
\[
P-P_0=(x-x_0,y-y_0,z-z_0).
\]
This displacement lies in the plane and is therefore perpendicular to
$\mathbf{n}$. Hence
\[
\mathbf{n}\cdot(P-P_0)=0.
\]

\begin{theorembox}
A plane through $P_0=(x_0,y_0,z_0)$ with normal vector
$\mathbf{n}=(a,b,c)\neq(0,0,0)$ has equation
\[
a(x-x_0)+b(y-y_0)+c(z-z_0)=0.
\]
After collecting constants, this becomes
\[
ax+by+cz=d
\]
for some $d\in\mathbb{R}$.
\end{theorembox}

Thus a first-degree equation in three variables describes a plane, just as a
first-degree equation in two variables describes a line.

\begin{examplebox}
The plane through $P_0=(1,-1,2)$ with normal vector
$\mathbf{n}=(2,1,-3)$ satisfies
\[
2(x-1)+(y+1)-3(z-2)=0.
\]
Therefore
\[
2x+y-3z=-5.
\]
\end{examplebox}

\subsection{A Plane Through Three Points}

Let $P$, $Q$, and $R$ be three points that do not lie on one line. Then
\[
\mathbf{u}=Q-P,
\qquad
\mathbf{v}=R-P
\]
are non-parallel direction vectors in the plane.

\begin{theorembox}
Three non-collinear points determine exactly one plane. A parametric equation is
\[
\Pi=\{P+s(Q-P)+t(R-P):s,t\in\mathbb{R}\}.
\]
A normal vector may be obtained from the cross product
\[
\mathbf{n}=(Q-P)\times(R-P).
\]
\end{theorembox}

\begin{examplebox}
For
\[
P=(1,0,0),
\qquad
Q=(0,1,0),
\qquad
R=(0,0,1),
\]
we have
\[
Q-P=(-1,1,0),
\qquad
R-P=(-1,0,1).
\]
Their cross product is
\[
(Q-P)\times(R-P)=(1,1,1).
\]
Using the point $P$, the plane equation is
\[
(x-1)+y+z=0,
\]
so
\[
x+y+z=1.
\]
\end{examplebox}

\subsection{Relative Positions of Planes}

Consider
\[
\Pi_1:a_1x+b_1y+c_1z=d_1,
\qquad
\Pi_2:a_2x+b_2y+c_2z=d_2,
\]
with normal vectors
\[
\mathbf{n}_1=(a_1,b_1,c_1),
\qquad
\mathbf{n}_2=(a_2,b_2,c_2).
\]

\begin{theorembox}
Two distinct planes are parallel precisely when their normal vectors are
parallel. If the normal vectors are not parallel, the planes intersect in a
line.
\end{theorembox}

The planes are perpendicular when their normal vectors are perpendicular:
\[
\mathbf{n}_1\cdot\mathbf{n}_2=0.
\]

\subsection{A Line as the Intersection of Two Planes}

In $\mathbb{R}^3$, two non-parallel planes usually meet in a line. Thus a line
may also be described by a system
\[
\begin{cases}
a_1x+b_1y+c_1z=d_1,\\
a_2x+b_2y+c_2z=d_2.
\end{cases}
\]
A direction vector for the intersection line is perpendicular to both normal
vectors, so one may take
\[
\mathbf{v}=\mathbf{n}_1\times\mathbf{n}_2.
\]

\begin{summarybox}
A plane in $\mathbb{R}^3$ may be described in two fundamental ways:
\[
P=P_0+s\mathbf{u}+t\mathbf{v}
\]
with two non-parallel direction vectors, or
\[
ax+by+cz=d
\]
with a nonzero normal vector $(a,b,c)$. Three non-collinear points determine one
plane, and two non-parallel planes intersect in a line.
\end{summarybox}
